\documentclass[12pt]{article}

\usepackage[utf8]{inputenc}
\usepackage[T1]{fontenc}
\usepackage{lmodern}
\usepackage[margin=1in]{geometry}
\usepackage{hyperref}
\usepackage{longtable}
\usepackage{array}
\usepackage{booktabs}

\renewcommand{\arraystretch}{1.2}
\newcommand{\toprule}{\hline}
\newcommand{\midrule}{\hline}
\newcommand{\bottomrule}{\hline}

\title{Security Plan: Pothole Detection System\\
\large Raspberry Pi Pipeline, Cloud Backend, and Frontend Applications}
\author{Capstone Project Team}
\date{July 2026}

\begin{document}
\maketitle

\begin{abstract}
This security plan identifies threats across the Raspberry~Pi vision pipeline,
the Vercel/Supabase cloud backend, and the Streamlit/React frontends; prioritizes
risks by likelihood and impact; defines authentication, authorization, input
validation, and data-security controls; and documents testing techniques, tools,
findings, mitigations, accepted risks, and re-test closure criteria aligned to
the course marking rubric.
\end{abstract}

\tableofcontents
\newpage

\section{Scope and System Boundary}

This report covers three cooperating parts of the pothole detection system:

\begin{enumerate}
    \item the Raspberry Pi vision-pipeline code presently in \texttt{rpi\_model/}
          and the planned SIM/cellular upload function;
    \item the cloud backend: a Vercel Next.js API (\texttt{api/}) backed by
          Supabase (Postgres table \texttt{pothole\_reports} and Storage bucket
          \texttt{pothole-images}); and
    \item the frontend applications that talk to that API: the Streamlit YOLO
          detection app (\texttt{backend/}) and the React map dashboard
          (\texttt{dashboard/}).
\end{enumerate}

Sensitive assets in scope are camera frames, annotated frames, pothole
detections, location data, model files, device and API credentials, Supabase
service secrets, map imagery URLs, and service availability.

\subsection{Data Flow}

\begin{verbatim}
Camera -> Raspberry Pi processing -> local Flask dashboard
                                   -> [planned] SIM/cellular HTTPS upload
                                                  -> Vercel API -> Supabase

Uploaded media -> Streamlit YOLO -> HTTPS POST (write key)
                                          -> Vercel API -> Supabase

React dashboard <- HTTPS GET (read key) <- Vercel API <- Supabase
\end{verbatim}

Hardware enclosure and physical deployment remain a placeholder for the hardware
engineer (Section~\ref{sec:hardware}).

\section{Threat Identification and Risk Prioritization}
\label{sec:threats}

\subsection{Threat Identification Method}

Threats were identified by reviewing attack surfaces on each trust boundary:
local network exposure of the Pi Flask dashboard, outbound cellular upload,
public HTTPS API routes, browser-origin calls to the API, uploaded image/
metadata payloads, and direct Supabase client access. Candidate threats include
missing authentication, privilege abuse, injection, insecure transport, secret
leakage, denial of service, and supply-chain compromise.

\subsection{Risk Prioritization Method}

Each threat is scored as Low, Medium, or High for \textbf{likelihood} and
\textbf{impact}. \textbf{Priority} is derived as follows:

\begin{itemize}
    \item \textbf{Critical}: High likelihood and High impact, or any path that
          yields unauthenticated write access to production data/storage;
    \item \textbf{High}: High impact with at least Medium likelihood, or High
          likelihood with Medium impact;
    \item \textbf{Medium}: Medium/Medium or lower combinations that still require
          tracked remediation or accepted-risk justification.
\end{itemize}

Every risk includes a \textbf{testing focus} so verification maps directly to the
threat.

\subsection{Pi and Cellular Upload Risk Register}

\begin{longtable}{p{0.22\textwidth}p{0.10\textwidth}p{0.10\textwidth}p{0.12\textwidth}p{0.36\textwidth}}
\caption{Pi and cellular-upload risk register}\label{tab:risks}\\
\toprule
Risk / threat & Likelihood & Impact & Priority & Testing focus \\
\midrule
\endfirsthead
\toprule
Risk / threat & Likelihood & Impact & Priority & Testing focus \\
\midrule
\endhead
Unauthenticated users can view the Pi dashboard, raw/annotated streams, and detection log. & High & High & Critical & Request every Flask route without credentials from a separate network client. \\
Camera frames and detection/location data travel over HTTP rather than HTTPS. & High & High & Critical & Capture traffic; confirm HTTP is refused in deployment. \\
Many concurrent MJPEG streams exhaust Pi CPU, memory, or bandwidth. & High & High & High & Concurrent-connection and bandwidth tests; verify connection limits. \\
Plaintext location lookup (\texttt{ip-api.com}) is untrusted and injectable into the UI. & Medium & High & High & Simulate modified, malformed, slow, and unavailable responses; XSS payload test. \\
Planned cellular uploader accepts forged, replayed, or modified requests. & Medium & High & High & Invalid credential, altered body, reused timestamp/nonce/request-ID tests. \\
Lost or copied device credentials submit false detections. & Medium & High & High & Per-device revocation and replacement test. \\
Unbounded configuration or malformed local inputs crash or starve the pipeline. & Medium & Medium & Medium & Boundary tests for dimensions, JPEG quality, inference settings, paths. \\
Unpinned dependencies or untrusted model checkpoints enable supply-chain compromise. & Medium & High & High & Dependency audit and model checksum/provenance verification. \\
Saved output frames fill local storage when \texttt{SAVE\_OUTPUT} is enabled. & Medium & Medium & Medium & Retention/quota and full-disk recovery tests. \\
\bottomrule
\end{longtable}

\subsection{Cloud Backend and Frontend Risk Register}

\begin{longtable}{p{0.22\textwidth}p{0.10\textwidth}p{0.10\textwidth}p{0.12\textwidth}p{0.36\textwidth}}
\caption{Cloud backend and frontend risk register}\label{tab:cloud-risks}\\
\toprule
Risk / threat & Likelihood & Impact & Priority & Testing focus \\
\midrule
\endfirsthead
\toprule
Risk / threat & Likelihood & Impact & Priority & Testing focus \\
\midrule
\endhead
Unauthenticated clients can \texttt{POST}/\texttt{PATCH}/\texttt{DELETE} reports and upload images. & High & High & Critical & Mutation routes without \texttt{X-API-Key}; expect 401. \\
Unauthenticated clients can enumerate reports via \texttt{GET /api/reports}. & High & Medium & High & GET without read key; expect 401 when \texttt{API\_READ\_KEY} is set. \\
Any website can call the API from a browser (CORS \texttt{*}). & High & High & High & OPTIONS/GET from an unlisted Origin; expect denial. \\
Upload endpoints accept spoofed \texttt{Content-Type} or oversized payloads. & Medium & High & High & Non-image bytes labelled JPEG; uploads $>$10\,MB. \\
Empty/fabricated detections become map markers and storage clutter. & Medium & Medium & High & POST with \texttt{detection\_count=0}; expect 400. \\
Supabase table readable/writable with a leaked anon key if RLS is off. & Medium & High & High & Query table with anon key after RLS enablement. \\
Write key or Supabase secret appears in frontend bundles or git. & Medium & Critical & Critical & Secret scan; confirm dashboard ships only the read key. \\
API abuse exhausts bandwidth, storage, or function quotas. & High & High & High & Burst requests; verify HTTP 429. \\
Production 500 responses leak internal exception text. & Medium & Medium & Medium & Forced server error; expect generic message. \\
XSS via unsanitized dashboard rendering of API/location text. & Medium & High & High & Inject script payloads into notes/location fields; verify text-only render. \\
SQL injection via report filters, IDs, or metadata if raw SQL were used. & Low & High & Medium & Parameterized ORM/client checks; malicious ID/query strings. \\
Command injection via filenames or configuration reaching a shell. & Low & High & Medium & Special-character filenames and env values; confirm no shell execution. \\
\bottomrule
\end{longtable}

\section{Authentication and Authorization}
\label{sec:auth}

\subsection{Current Findings: Weak or Missing Authentication}

\paragraph{Raspberry Pi Flask dashboard.}
The Pi application currently has \textbf{no authentication}, passwords, tokens,
roles, or route-level authorization. It binds to \texttt{0.0.0.0}, exposing:

\begin{itemize}
    \item \texttt{/} --- dashboard;
    \item \texttt{/stream/original.mjpg} --- raw camera feed;
    \item \texttt{/stream/edges.mjpg} --- annotated camera feed;
    \item \texttt{/detections.json} --- recent detections and location text.
\end{itemize}

This is a \textbf{Critical} finding if the Pi is reachable beyond a trusted
maintenance network. Until authentication is added, access must be limited to
localhost or a firewall/VPN-protected network.

\paragraph{Cloud API (implemented).}
The Vercel API previously accepted unauthenticated writes. It now enforces
API-key authentication:

\begin{itemize}
    \item \textbf{Write privilege}: \texttt{API\_WRITE\_KEY} required for
          \texttt{POST}/\texttt{PATCH}/\texttt{DELETE} and
          \texttt{/api/reports/backfill-locations};
    \item \textbf{Read privilege}: \texttt{API\_READ\_KEY} required for
          \texttt{GET /api/reports} and \texttt{GET /api/reports/[id]} when
          configured (enforced in production);
    \item \textbf{Health}: \texttt{GET /api/health} remains unauthenticated but
          rate-limited.
\end{itemize}

Keys are compared with constant-time equality. Clients send
\texttt{X-API-Key} or \texttt{Authorization: Bearer}.

\subsection{Access Privileges}

\begin{longtable}{p{0.24\textwidth}p{0.28\textwidth}p{0.38\textwidth}}
\caption{Access privilege matrix}\label{tab:privileges}\\
\toprule
Principal & Credential / channel & Privileges \\
\midrule
\endfirsthead
\toprule
Principal & Credential / channel & Privileges \\
\midrule
\endhead
Anonymous network client (Pi) & None & Full local dashboard/stream access today (unacceptable for public networks). \\
Streamlit detection app & \texttt{POTHOLE\_API\_WRITE\_KEY} & Create reports only (POST multipart). \\
React dashboard / browser & \texttt{VITE\_API\_READ\_KEY} & Read reports only (GET). No mutate. \\
Vercel API process & \texttt{SUPABASE\_SECRET\_KEY} & Server-side insert/update/delete/storage; never shipped to browsers. \\
Browser with Supabase anon key & Publishable/anon key & No table access when RLS is enabled with no public policies. \\
Future Pi cellular device & Per-device revocable credential & Upload only to fixed HTTPS ingestion URL; no inbound server role. \\
\bottomrule
\end{longtable}

\subsection{Password Entropy}

\begin{itemize}
    \item \textbf{Pi dashboard today:} password entropy is \emph{not applicable}
          because no password login exists. If a local password is introduced, it
          must be hashed with Argon2id or bcrypt, rate-limited, never stored in
          source, and meet a minimum entropy policy (at least 12 characters with
          mixed classes, or a $\geq$128-bit random secret).
    \item \textbf{Cloud API keys:} not user passwords; they are long random
          secrets (\texttt{secrets.token\_urlsafe}, $\geq$24--32 bytes). Entropy
          derives from cryptographic randomness rather than human-chosen
          passwords. Keys are rotated by regenerating and updating Vercel/local
          env vars.
    \item \textbf{Future end-user accounts:} if Supabase Auth/OAuth is added,
          enforce provider password policy or SSO, MFA for operators, and hashed
          storage managed by the auth provider.
\end{itemize}

\subsection{Authentication and Authorization Test Cases}

\begin{longtable}{p{0.08\textwidth}p{0.42\textwidth}p{0.40\textwidth}}
\caption{Authentication and authorization tests}\label{tab:auth-tests}\\
\toprule
ID & Test & Expected result \\
\midrule
\endfirsthead
\toprule
ID & Test & Expected result \\
\midrule
\endhead
A-01 & Call every Pi Flask route with no credentials from another host. & Today: access succeeds (documented vulnerability). Target: denial after auth is added. \\
A-02 & \texttt{GET /api/reports} with no key. & HTTP 401. \\
A-03 & \texttt{POST /api/reports} with no key. & HTTP 401. \\
A-04 & \texttt{GET /api/reports} with valid read key. & HTTP 200. \\
A-05 & \texttt{POST /api/reports} with read key only. & HTTP 401 (read cannot elevate to write). \\
A-06 & \texttt{DELETE /api/reports/[id]} with write key missing/wrong. & HTTP 401. \\
A-07 & Dashboard bundle/source review for write/secret keys. & Only read key present; no Supabase secret. \\
A-08 & Streamlit post with write key unset. & Client refuses before network call. \\
A-09 & (Future) weak Pi password / brute-force attempts. & Lockout/rate-limit; hashes only stored. \\
\bottomrule
\end{longtable}

\textbf{Evidence already collected for cloud auth:} production
\texttt{https://api-mu-ten-54.vercel.app} returns 401 for unauthenticated GET/POST
and 200 for authorized GET with the read key.

\section{Input Validation and Injection Testing}
\label{sec:validation}

\subsection{Validation Strategy: Client-Side and Server-Side}

Security-critical validation is enforced on the \textbf{server} (Vercel API /
future cloud ingestion). Clients may pre-check for usability, but are never
trusted:

\begin{itemize}
    \item \textbf{Streamlit (client):} only posts when detections contain
          ``pothole''; refuses to post without a write key; trims notes.
    \item \textbf{React dashboard (client):} filters markers to records with
          coordinates and \texttt{detectionCount > 0}; does not accept free-form
          write inputs to the API.
    \item \textbf{API (server):} validates content type, image magic bytes, size
          ($\leq$10\,MB), metadata JSON shape/size, severity/status enums, notes
          length ($\leq$2000), and rejects empty detections with HTTP~400.
    \item \textbf{Pi (local):} must add range checks for env-derived numeric
          config (dimensions, JPEG quality 1--100, confidence 0--1, etc.).
\end{itemize}

\subsection{Injection Attack Surfaces}

\begin{itemize}
    \item \textbf{SQL injection.} The API uses the Supabase client with
          parameterized queries/filters (no string-concatenated SQL). Direct SQL
          from browsers is blocked by RLS. Regression tests still send malicious
          IDs and query strings.
    \item \textbf{Cross-site scripting (XSS).} Pi dashboard historically used
          \texttt{innerHTML} for external location text --- high risk. Cloud
          notes/metadata must be rendered as text in the React UI. CORS
          allowlisting reduces hostile-site scripted API abuse from browsers.
    \item \textbf{Command injection.} Neither Pi nor API currently shells out on
          user input; filenames are sanitized. Tests still use
          \texttt{; rm -rf}, backticks, and newline names to confirm no shell
          path exists.
    \item \textbf{Other input abuse.} Oversized JSON, type confusion
          (array instead of object metadata), spoofed MIME types, path traversal
          in filenames.
\end{itemize}

\subsection{Input Validation and Injection Test Cases}

\begin{longtable}{p{0.08\textwidth}p{0.42\textwidth}p{0.40\textwidth}}
\caption{Input validation and injection tests}\label{tab:inject-tests}\\
\toprule
ID & Test vector & Expected result \\
\midrule
\endfirsthead
\toprule
ID & Test vector & Expected result \\
\midrule
\endhead
I-01 & POST metadata with \texttt{detection\_count=0}. & HTTP 400; no row created. \\
I-02 & Upload \texttt{.exe}/HTML bytes with \texttt{Content-Type: image/jpeg}. & Rejected (magic-byte mismatch). \\
I-03 & Upload image $>$10\,MB. & Rejected. \\
I-04 & Metadata JSON array or oversized blob. & HTTP 400. \\
I-05 & Invalid severity/status strings. & HTTP 400. \\
I-06 & Notes with \texttt{<script>alert(1)</script>}. & Stored only if authorized write; dashboard renders as text, does not execute. \\
I-07 & Pi location response containing script markup. & Must display via \texttt{textContent}; script must not run. \\
I-08 & Report ID \texttt{1; DROP TABLE pothole\_reports;--}. & No SQL execution; 404/validation error only. \\
I-09 & Filename \texttt{../../etc/passwd;id.jpg}. & Sanitized object key; no path escape / no shell. \\
I-10 & Extreme Pi env values (negative sizes, quality 999). & Safe rejection or clamped config (target after PI-07). \\
\bottomrule
\end{longtable}

\section{Data Security and Privacy Checks}
\label{sec:data}

\subsection{Secure Transmission (HTTPS / TLS)}

\begin{itemize}
    \item \textbf{Cloud API and frontends:} production traffic uses HTTPS
          (\texttt{https://api-mu-ten-54.vercel.app}). Streamlit posts and
          dashboard fetches must target HTTPS URLs; no HTTP fallback for cloud
          calls.
    \item \textbf{Pi today:} local Flask dashboard and
          \texttt{http://ip-api.com} location lookup use plaintext HTTP --- a
          Critical data-exposure risk on untrusted networks.
    \item \textbf{Planned SIM upload:} HTTPS only, with certificate validation;
          HTTP forbidden.
\end{itemize}

\subsection{Encryption in Transit and at Rest}

\begin{itemize}
    \item \textbf{In transit:} TLS terminated by Vercel and Supabase for API and
          database/storage connections from the server.
    \item \textbf{At rest:} managed encryption at rest provided by Supabase /
          underlying cloud storage for Postgres data and bucket objects.
    \item \textbf{Secrets at rest on developer machines:} stored in gitignored
          \texttt{.env} / \texttt{.env.local} files and Vercel encrypted
          environment variables---not in the repository.
\end{itemize}

\subsection{Data Leakage and Sensitive Information Exposure}

Sensitive classes include camera imagery (people, plates, addresses), GPS /
temporary location, API keys, and Supabase secrets.

Controls and checks:

\begin{itemize}
    \item Frontends never embed \texttt{SUPABASE\_SECRET\_KEY} or
          \texttt{API\_WRITE\_KEY}.
    \item Production API errors return generic messages (no stack traces or
          provider internals).
    \item Rate limits and auth reduce anonymous bulk export of reports.
    \item Pi must not expose streams on public/cellular networks until auth and
          TLS exist; minimize retained frames and define retention before
          enabling \texttt{SAVE\_OUTPUT}.
    \item Public storage URLs are an accepted demo risk with compensating
          write-path controls; private signed URLs are the longer-term fix.
\end{itemize}

\subsection{Data Security Test Cases}

\begin{longtable}{p{0.08\textwidth}p{0.42\textwidth}p{0.40\textwidth}}
\caption{Data security and privacy tests}\label{tab:data-tests}\\
\toprule
ID & Test & Expected result \\
\midrule
\endfirsthead
\toprule
ID & Test & Expected result \\
\midrule
\endhead
D-01 & Confirm production API URL scheme is HTTPS. & \texttt{https://} only. \\
D-02 & Attempt HTTP cleartext to production API. & Redirect or connection failure; no successful cleartext API use. \\
D-03 & Capture Streamlit upload and dashboard fetch. & TLS; no write key in dashboard traffic. \\
D-04 & Force API 500 in production. & Generic error body; no secret/stack leakage. \\
D-05 & Search git history / bundles for secret patterns (Gitleaks). & No live secrets committed. \\
D-06 & Query Supabase with anon key after RLS. & Access denied to \texttt{pothole\_reports}. \\
D-07 & Pi traffic capture on LAN. & Today: plaintext imagery/location visible (open finding PI-02). \\
\bottomrule
\end{longtable}

\section{Testing Techniques and Tools}
\label{sec:testing}

This section describes \emph{how} the threats in Section~\ref{sec:threats} are
verified. Techniques are applied across Pi, API, and frontends.

\subsection{Static Application Security Testing (SAST)}

SAST analyzes source without executing it. Planned/used tools:

\begin{itemize}
    \item \textbf{Bandit} --- Python SAST for the Pi and Streamlit code (unsafe
          functions, hard-coded secrets, weak crypto patterns).
    \item \textbf{Semgrep} --- rule-based SAST for Python/TypeScript (Flask
          security, insecure HTTP, dangerous DOM sinks).
    \item \textbf{ESLint security plugins / TypeScript compiler} --- dashboard
          and API static checks.
    \item \textbf{Ruff} --- already used for Pi linting; complements SAST hygiene.
\end{itemize}

Evidence: clean reports or documented false positives / accepted risks.

\subsection{Dynamic Application Security Testing (DAST)}

DAST exercises a running service with crafted HTTP requests:

\begin{itemize}
    \item \textbf{OWASP ZAP} --- automated spider/active scan against the Pi
          Flask dashboard (when safely reachable on a lab network) and against a
          staging copy of the cloud API.
    \item \textbf{Burp Suite} --- manual intercept/replay of API requests to
          tamper with headers, bodies, cookies/tokens, and CORS preflights.
    \item \textbf{curl / HTTPie scripts} --- repeatable regression checks for
          401/403/400/429 behaviour on production and local API.
\end{itemize}

DAST focuses on missing authentication, privilege escalation (read key used as
write), injection payloads, TLS downgrade attempts, and information leakage in
responses.

\subsection{Manual Code Review}

Reviewers inspect Flask routes, Next.js API handlers (\texttt{api/lib/auth.ts},
\texttt{storage.ts}, \texttt{http.ts}), Streamlit upload client, React fetch
paths, logging, model loading, and systemd/service configuration. Checklist
items: secrets in source, auth on every mutating route, validation before
persistence, safe rendering, and privilege separation.

\subsection{Input Validation Checks}

Covered in Section~\ref{sec:validation}: boundary values, type confusion,
magic-byte mismatches, and injection strings are sent to both client pre-checks
(expect soft failure) and server endpoints (expect hard rejection).

\subsection{Authentication and Authorization Testing}

Covered in Section~\ref{sec:auth}: missing credentials, wrong credentials, wrong
privilege class, and future password-entropy/brute-force cases.

\subsection{Additional Supporting Tools}

\begin{longtable}{p{0.28\textwidth}p{0.32\textwidth}p{0.30\textwidth}}
\caption{Security tools and required evidence}\label{tab:tools}\\
\toprule
Tool / technique & Application & Required evidence \\
\midrule
\endfirsthead
\toprule
Tool / technique & Application & Required evidence \\
\midrule
\endhead
Manual code review & Pi, API, Streamlit, dashboard & Review notes linked to finding IDs. \\
Bandit / Semgrep (SAST) & Python/TS sources & Scan report or waived findings. \\
OWASP ZAP (DAST) & Pi dashboard; staging API & Scan summary + false-positive review. \\
Burp Suite (DAST) & API request tampering & Request/response exports for A-/I-/D-cases. \\
\texttt{pip-audit} / \texttt{npm audit} & Dependencies & Audit output; pinned versions. \\
Gitleaks & Repo + history & No unreviewed secrets. \\
curl regression scripts & Production API auth/CORS/validation & Timestamped command output. \\
Load / concurrency tests & Pi MJPEG; API rate limits & Resource metrics; HTTP 429 proof. \\
\bottomrule
\end{longtable}

\section{Subsystem Detail: Raspberry Pi and Planned SIM Upload}
\label{sec:pi}

\subsection{Pi Input Validation and Injection (Current Code)}

Numeric environment values are converted with \texttt{int()}/\texttt{float()}
without range checks. Extreme values can exhaust resources. The location lookup
response is inserted with JavaScript \texttt{innerHTML}, enabling XSS if the
HTTP response is tampered with. There are presently no SQL or shell sinks on
user input; those remain regression requirements if such features appear.

\subsection{Pi Data Security and Availability}

Plaintext HTTP for the dashboard and location service exposes imagery and
location. Frames may contain PII; retention must be minimized. Uncapped MJPEG
viewers create denial-of-service risk. systemd hardening should use a non-root
account, \texttt{NoNewPrivileges=true}, and narrowly writable output paths.

\subsection{Planned SIM/Cellular Upload Requirements}

The SIM module is outbound-only to a fixed HTTPS cloud URL. Each Pi needs a
unique revocable credential (not a shared fleet key), TLS, timestamp + nonce /
request ID, body integrity (e.g.\ signed hash), and rotation/revocation. The
cloud must reject unknown devices, bad signatures, stale timestamps, duplicate
nonces, and oversized/invalid payloads. Queuing uses bounded retries and must
never log credentials.

\section{Subsystem Detail: Backend and Cloud Data Service}
\label{sec:backend}
\textit{Owner: Backend engineer (Amir Abdalla).}

Production endpoint: \texttt{https://api-mu-ten-54.vercel.app}. Source:
\texttt{api/}.

\subsection{Implemented Controls}

\begin{enumerate}
    \item Least privilege: browsers never receive the Supabase secret.
    \item Separate read/write API keys with fail-closed writes.
    \item CORS Origin allowlist (local Vite defaults +
          \texttt{CORS\_ALLOWED\_ORIGINS}).
    \item Per-IP rate limiting (stricter on writes); HTTP 429 +
          \texttt{Retry-After}.
    \item Upload hardening: JPEG/PNG/WEBP, 10\,MB max, magic-byte verification,
          sanitized UUID object keys.
    \item Schema checks: metadata caps, enums, notes length, empty-detection
          rejection.
    \item Generic production 500 responses.
    \item RLS script for \texttt{pothole\_reports}
          (\texttt{api/supabase/rls\_pothole\_reports.sql}).
\end{enumerate}

\subsection{Backend Verification Evidence}

\begin{itemize}
    \item Unauthenticated GET/POST on production $\rightarrow$ HTTP 401.
    \item Authorized GET with read key $\rightarrow$ HTTP 200.
    \item Write key + empty detections $\rightarrow$ HTTP 400.
    \item CORS preflight from unlisted Origin $\rightarrow$ HTTP 403.
    \item CORS preflight from \texttt{http://localhost:5173} $\rightarrow$
          HTTP 204 with allowlisted ACAO header.
\end{itemize}

\section{Subsystem Detail: Frontend Integration}
\label{sec:frontend}
\textit{Owner: Frontend / detection-app engineers.}

\begin{itemize}
    \item Streamlit loads \texttt{POTHOLE\_API\_WRITE\_KEY} from
          \texttt{backend/.env} (never shown in UI) and posts over HTTPS.
    \item Dashboard uses \texttt{VITE\_API\_READ\_KEY} only; GET-only; no
          Supabase admin client.
    \item Client-side filters complement, but do not replace, server validation.
    \item Accepted limitation: the Vite read key is visible in the browser
          bundle and is treated as a shared team token, not a root secret
          (see AR-01).
\end{itemize}

\section{Reporting and Risk Mitigation Plan}
\label{sec:reporting}

\subsection{Vulnerability and Finding Register}

Findings are tracked with severity, status, owner, and closure evidence. Status
values: \textbf{Open}, \textbf{Mitigated}, \textbf{Accepted} (with
justification), or \textbf{Closed} (fix + re-test passed).

\begin{longtable}{p{0.10\textwidth}p{0.10\textwidth}p{0.10\textwidth}p{0.32\textwidth}p{0.28\textwidth}}
\caption{Pi/SIM findings and mitigations}\label{tab:mitigations}\\
\toprule
ID & Severity & Status & Required mitigation & Closure / re-test evidence \\
\midrule
\endfirsthead
\toprule
ID & Severity & Status & Required mitigation & Closure / re-test evidence \\
\midrule
\endhead
PI-01 & Critical & Open & Restrict dashboard network; add authentication/authorization. & Unauthenticated route tests denied; proxy/firewall reviewed. \\
PI-02 & Critical & Open & HTTPS for all remote connections; replace plaintext location lookup. & Traffic capture shows TLS only. \\
PI-03 & High & Open & Render location with DOM text APIs; validate response schema. & XSS payload displays as text only. \\
PI-04 & High & Open & Cap MJPEG viewers; monitor resources. & Load test within agreed limits. \\
PI-05 & High & Open & Per-device SIM auth, integrity, replay protection, revocation. & Forged/stale/duplicate uploads rejected. \\
PI-06 & High & Open & Pin/audit dependencies; verify model checksums. & Manifest + audit + checksum log. \\
PI-07 & Medium & Open & Validate configuration ranges; safe failure messages. & Automated boundary tests pass. \\
PI-08 & Medium & Open & Output retention/quota; least-privilege service account. & Full-disk/retention tests; permission review. \\
\bottomrule
\end{longtable}

\begin{longtable}{p{0.10\textwidth}p{0.10\textwidth}p{0.12\textwidth}p{0.30\textwidth}p{0.28\textwidth}}
\caption{Cloud/frontend findings and mitigations}\label{tab:cloud-mitigations}\\
\toprule
ID & Severity & Status & Required mitigation & Closure / re-test evidence \\
\midrule
\endfirsthead
\toprule
ID & Severity & Status & Required mitigation & Closure / re-test evidence \\
\midrule
\endhead
BE-01 & Critical & Closed & Write API key on mutations; fail closed. & Prod POST without key $\rightarrow$ 401. \\
BE-02 & High & Closed & Read API key on report GETs. & Prod GET without key $\rightarrow$ 401; with key $\rightarrow$ 200. \\
BE-03 & High & Closed & CORS allowlist instead of \texttt{*}. & Unlisted Origin $\rightarrow$ 403; local Vite allowed. \\
BE-04 & High & Mitigated & Per-IP rate limits on API routes. & Burst returns 429 (in-memory limiter). \\
BE-05 & High & Closed & Magic-byte + type + size upload checks. & Spoofed/oversized uploads rejected. \\
BE-06 & High & Mitigated & Enable Supabase RLS; secret key server-only. & SQL script provided; confirm applied in project. \\
BE-07 & Medium & Closed & Cap metadata/notes; reject empty detections. & Empty detection POST $\rightarrow$ 400. \\
BE-08 & Medium & Closed & Generic production 500 bodies. & No stack/provider leakage in prod errors. \\
FE-01 & High & Closed & Streamlit write key from env only. & Missing key blocks post; secret not in UI. \\
FE-02 & High & Closed & Dashboard read-only API use. & Code review + GET-only traffic. \\
FE-03 & Medium & Accepted & Document bundlable read key as team token. & See AR-01. \\
\bottomrule
\end{longtable}

\subsection{Accepted Risks (with Justification)}

\begin{longtable}{p{0.10\textwidth}p{0.28\textwidth}p{0.28\textwidth}p{0.24\textwidth}}
\caption{Accepted risks}\label{tab:accepted}\\
\toprule
ID & Risk & Justification / compensating control & Review date \\
\midrule
\endfirsthead
\toprule
ID & Risk & Justification / compensating control & Review date \\
\midrule
\endhead
AR-01 & Dashboard read key visible in JS bundle. & Mid-project shared demo; key is read-only; writes still fail closed; rotate if leaked. & End of Week 12 \\
AR-02 & Public Supabase image URLs. & Writes API-gated; RLS on table; migrate to signed URLs before public launch. & Before public demo \\
AR-03 & No end-user login / RBAC yet. & Operator-only demo; API keys separate read/write; track Auth as future work. & Final delivery \\
AR-04 & In-memory API rate limiter. & Adequate for single-region demo; distributed limiter/WAF later. & Production hardening \\
\bottomrule
\end{longtable}

\subsection{Re-testing and Issue Closure Process}

An issue is \textbf{Closed} only when all of the following are true:

\begin{enumerate}
    \item The fix is implemented and code-reviewed.
    \item The mapped test IDs (A-/I-/D- or table testing focus) are re-run.
    \item Evidence (request/response, scan excerpt, or review note) is attached to
          the finding ID.
    \item No regression appears in subsequent CI/manual smoke tests of auth,
          validation, and TLS.
\end{enumerate}

\textbf{Accepted} risks require an owner, justification, compensating control,
and a scheduled review date (Table~\ref{tab:accepted}). Open Pi items (PI-01--PI-08)
remain the active remediation backlog for the edge device and SIM path.

\section{Deferred Production Hardening}

\begin{itemize}
    \item End-user authentication (Supabase Auth / OAuth) and RBAC views.
    \item Per-device Pi credentials with nonce/replay protection on cellular
          upload.
    \item Private storage with short-lived signed URLs.
    \item Distributed rate limiting / WAF; CI-wired Bandit, Semgrep, ZAP, and
          Gitleaks.
\end{itemize}

\section{Hardware and Physical Deployment}
\label{sec:hardware}
\textit{Owner: \underline{\hspace{4cm}} (to be completed by the hardware engineer).}

This section must specify physical tamper controls, power protection,
camera/privacy placement, cellular SIM management, exposed-port controls,
lost-device response, and maintenance procedures.

\appendix
\section{Rubric Coverage Checklist}

\begin{longtable}{p{0.28\textwidth}p{0.62\textwidth}}
\caption{Mapping of this plan to the marking rubric}\label{tab:rubric}\\
\toprule
Rubric criterion & Where addressed \\
\midrule
\endfirsthead
\toprule
Rubric criterion & Where addressed \\
\midrule
\endhead
Threat identification & Section~\ref{sec:threats}; Tables~\ref{tab:risks}, \ref{tab:cloud-risks} \\
Risk prioritization (likelihood/impact) & Section~\ref{sec:threats} method + priority columns \\
Risk $\leftrightarrow$ testing focus mapping & Final column of both risk registers \\
SAST / DAST / manual / validation / auth testing & Section~\ref{sec:testing} \\
Named tools (ZAP, Burp, Bandit, \ldots) & Section~\ref{sec:testing}, Table~\ref{tab:tools} \\
Weak/missing auth tests & Section~\ref{sec:auth}, Table~\ref{tab:auth-tests} \\
Password entropy & Section~\ref{sec:auth} \\
Access privileges & Table~\ref{tab:privileges} \\
SQL injection / XSS / command injection & Section~\ref{sec:validation}, Table~\ref{tab:inject-tests} \\
Client- and server-side validation & Section~\ref{sec:validation} \\
HTTPS/TLS; encryption at rest/in transit & Section~\ref{sec:data} \\
Data leakage checks & Section~\ref{sec:data}, Table~\ref{tab:data-tests} \\
Findings with severity; fixes; accepted risks; re-test closure & Section~\ref{sec:reporting} \\
\bottomrule
\end{longtable}

\end{document}
